\chapter{指标定理、规范场与进阶专题}
\label{ch:20-index-gauge}
\ChapterPlan{周第10章（本地页码待核对）；Dai的Dirac讲义第3章、Figueroa-O'Farrill的2006规范场讲义3--5、Perutz的2018讲义第17--20节}{正文证明初稿；8题详解、3幅原创图；深层指标公式外引，独立数学审阅未完成}

第19章的Dirac算子有可逆主符号，但这并不自动说明它的像闭或核有限维。本章先把局部微分性质接到全局Hilbert空间问题，再说明指标公式所增加的内容。另一条线从联络仿射空间和特征形式出发，推导Yang--Mills方程、四维能量界及规范固定的线性符号。最后在给定Spin$^c$数据下建立Seiberg--Witten方程及消失结论，不宣称已经构造四维光滑不变量。

除显式反例外，$M$为紧致、光滑、无边界的黎曼流形；谈积分特征数及四维自对偶时另给定定向。复丛有Hermite内积，联络保持内积。沿用$\partial_j\mapsto i\xi_j$的主符号约定、$c(v)^2=-|v|^2$和$R=[\nabla,\nabla]-\nabla_{[\ ,\ ]}$。$H^s$为第14章的Sobolev空间；Lebesgue积分、$L^2$完备性仍是分析先修。

\section{椭圆算子、符号与Fredholm指标}
\subsection{主符号与需要选择的函数空间}
若局部$P=\sum_{|\alpha|\leq m}a_\alpha(x)\partial^\alpha$，定义
\[
 \sigma_m(P)(x,\xi)=\sum_{|\alpha|=m}a_\alpha(x)(i\xi)^\alpha.
\]
坐标变换的链式法则只有最高阶导数贡献于此式；换丛标架则在两端乘相应变换矩阵。故它是$\pi^*E\to\pi^*F$的丛同态，$\pi:T^*M\to M$。当每个$\xi\ne0$的符号可逆时称$P$椭圆（elliptic）；本章此定义要求两端同秩。相乘最高阶项给$\sigma(PQ)=\sigma(P)\sigma(Q)$，分部积分给$\sigma(P^*)=\sigma(P)^*$。

光滑截面空间尚未带本章所用的Hilbert范数。光滑系数及紧致性使每个$m$阶算子有有界延拓$P:H^{s+m}(E)\to H^s(F)$，$s\geq0$为整数：有限覆盖中应用Leibniz公式，每项至多是有界光滑系数乘$u$的不超过$s+m$阶弱导数。一般椭圆算子的完整参数子理论暂不建立；以下直接证明本书所需的一阶Dirac型情形。

设$P:E\to F$为一阶算子，满足
\begin{equation}\label{eq:index-dirac-type}
 \sigma_P(x,\xi)^*\sigma_P(x,\xi)=|\xi|^2I_E,\qquad
 \sigma_P(x,\xi)\sigma_P(x,\xi)^*=|\xi|^2I_F.
\end{equation}
第19章的$D^+$、$D^-$及其酉扭曲、加光滑零阶项均满足该式；$\dd+\delta$的奇偶分块也满足。$P^*$在此始终指形式伴随，不是已经指定定义域的无界算子伴随。

\Needspace{14\baselineskip}
\begin{lemma}[一阶估计及分布核的光滑性]\label{lem:index-first-order}
对满足\eqref{eq:index-dirac-type}的$P$，有
\[
 \|u\|_{H^1}\leq C(\|Pu\|_{L^2}+\|u\|_{L^2}),\qquad u\in H^1(E).
\]
若$u\in L^2_{\mathrm{loc}}$且$Pu\in L^2_{\mathrm{loc}}$（分布意义），则$u\in H^1_{\mathrm{loc}}$。特别地，分布核$Pu=0$中的$L^2_{\mathrm{loc}}$截面光滑。相同结论适用于$P^*$。
\end{lemma}
\begin{proof}
先对支撑在小坐标球的光滑$u$冻结一阶系数，所得$P_0=\sum A_j(x_0)\partial_j$满足
$\|P_0u\|_2^2=\int|\xi|_{g(x_0)}^2|\widehat u(\xi)|^2\dd\xi$。
此处可把支撑包入较大周期盒，使用第14章的Fourier论证。坐标与度量范数等价，缩小球后
\[
 \|(P-P_0)u\|_2\leq\epsilon\|\partial u\|_2+C_\epsilon\|u\|_2.
\]
取$\epsilon$足够小并吸收，得到局部估计。有限单位分解的交换子$[P,\chi]$为零阶，求和给全局估计，再由光滑稠密性延拓到$H^1$。

要应用于分布核，不能先假定$u\in H^1$。取内部截断$w=\chi u$并用卷积$J_\varepsilon w=\rho_\varepsilon*w$。对一个一阶项，分布积分给
\begin{align*}
 [A_j\partial_j,J_\varepsilon]w(x)
 ={}&\int(A_j(x)-A_j(y))\partial_j\rho_\varepsilon(x-y)w(y)\dd y\\
 &+\int\rho_\varepsilon(x-y)(\partial_jA_j)(y)w(y)\dd y.
\end{align*}
在较小坐标球延拓系数为有界光滑系数。第一核被$C|x-y||\partial\rho_\varepsilon(x-y)|$控制，其积分与$\varepsilon$无关；第二核积分也有统一界。Young卷积不等式（由Cauchy--Schwarz和Fubini得到）给交换子的统一$L^2$界，零阶项同样有界。因此
\[
 \|P J_\varepsilon w\|_2\leq C(\|Pw\|_2+\|w\|_2),\qquad
 Pw=\chi Pu+[P,\chi]u\in L^2.
\]
对平滑截面应用局部一阶估计，得统一$H^1$界。弱收敛子列的$L^2$极限是$w$，分布导数的极限是其弱导数，所以$w\in H^1$。

最后$Pu=0$蕴含$P^*Pu=0$。该二阶算子的主部为$-g^{ij}\partial_i\partial_j I$，可写成第14章引理\ref{lem:hodge-weak-regularity}的标量主部系统；其差商证明只使用该主部及光滑低阶系数，逐阶应用得到所有局部$H^s$正则性。Sobolev嵌入给光滑性。交换$P,P^*$完成另一端。
\end{proof}

零阶项不改主符号，但可改变核。例如单位长度圆周$P_a=\partial_x+ia$在周期复函数上的核，只有$a\in2\pi\mathbb Z$时非零。椭圆性约束高频，不能单独决定零频。非紧时还会失去紧嵌入，后面的Fredholm结论须重新论证。

\subsection{从紧嵌入到Fredholm}
Hilbert空间之间的有界算子$T:X\to Y$称为Fredholm算子，当它的像闭且核与余核$Y/\operatorname{im}T$均有限维；定义
\[
 \operatorname{ind}T=\dim\ker T-\dim(Y/\operatorname{im}T).
\]
闭像是定义的一部分，不能只写两个有限维条件而忽略所用拓扑。

\Needspace{16\baselineskip}
\begin{theorem}\label{thm:index-fredholm}
在本章紧致无边界假设下，满足\eqref{eq:index-dirac-type}的$P:H^1(E)\to L^2(F)$是Fredholm，且
\[
 (\operatorname{im}P)^{\perp_{L^2}}=\ker P^*,\qquad
 \operatorname{ind}P=\dim\ker P-\dim\ker P^*.
\]
两端的核均由光滑截面组成。
\end{theorem}
\begin{proof}
第14章的紧嵌入$H^1\hookrightarrow L^2$适用于任意有限秩丛：有限平凡化把它化为有限个标量结论。核上一阶估计给$\|u\|_{H^1}\leq C\|u\|_2$。若核无限维，取$L^2$规范正交序列；它在$H^1$有界，却无$L^2$收敛子列，矛盾。所以$K=\ker P$有限维，且由引理\ref{lem:index-first-order}光滑。

在$K$的$L^2$正交补上有更强估计
\begin{equation}\label{eq:index-gap}
 \|u\|_{H^1}\leq C\|Pu\|_2,\qquad u\perp K.
\end{equation}
否则取$\|u_j\|_{H^1}=1$、$u_j\perp K$及$Pu_j\to0$。紧嵌入给$L^2$收敛子列，一阶估计用于$u_j-u_k$使其在$H^1$中Cauchy。极限$u$属于$K\cap K^\perp$且$H^1$范数1，矛盾。

若$Pu_j\to f$，先减去$u_j$到$K$的有限维投影；\eqref{eq:index-gap}使余下序列在$H^1$收敛到$u$，连续性给$Pu=f$，故像闭。若$v\in L^2$正交于像，则对所有光滑$u$有$(Pu,v)=0$，即$P^*v=0$在分布意义成立。上个引理使$v$光滑，故它是$P^*$的核；反向由分部积分成立。对$P^*$重复有限维核的论证，余核有限维。
\end{proof}

这一证明没有把形式自伴直接等同于本质自伴，也未证明完整谱展开。偶维手征算子得
\[
 \operatorname{ind}D_E^+=\dim\ker D_E^+-\dim\ker D_E^-.
\]
全$D_E$的指标为零，手征指标却未必为零。对$\dd+\delta$的偶到奇分块，第14章Hodge分解独立给出指标$\sum_k(-1)^k\dim H^k_{\mathrm{dR}}(M)$；这是解析指标的计算，尚未证明其等于Euler形式积分。

\begin{figure}[htbp]
 \centering\DGDiagram{fredholm-splitting}

 \caption{Fredholm算子的有限维缺陷（原创示意）。中间限制为有界同构，核与余核的维数之差才是指标；图中方块不表示相等维数。}
\end{figure}

\subsection{指标为何对扰动稳定}
\begin{proposition}\label{prop:index-stability}
Fredholm算子在算子范数的小扰动下仍Fredholm且指标不变。加紧算子保持Fredholm及指标。因此固定丛、固定Sobolev空间之间的连续Dirac型算子族具有常指标；特别地，更改酉扭联络或加光滑零阶项不改变指标。
\end{proposition}
\begin{proof}
写$X=K\oplus X_0$、$Y=C\oplus Y_0$，其中$Y_0=\operatorname{im}T$、$C=Y_0^\perp$。$T_0:X_0\to Y_0$为有界同构。在本章情形逆有界来自\eqref{eq:index-gap}；一般情形用有界逆定理。此定理是Hilbert分析工具，可由Baire定理证明，不是有限维逆函数定理；本章一般Hilbert版本将它列为外部分析先修。

对$S$充分接近$T$，其$X_0\to Y_0$块$d$仍可逆，因为$T_0^{-1}d=I+R$且$\|R\|<1$，逆由收敛Neumann级数给出。写$S=\left(\begin{smallmatrix}a&b\\c&d\end{smallmatrix}\right)$，两侧乘有界可逆三角矩阵消元，得到对角块$a-bd^{-1}c:K\to C$及$d:X_0\to Y_0$。有限维第一块的指标恒为$\dim K-\dim C$，第二块指标零；也得到闭像和有限维缺陷。这证明局部稳定性。

定义$B:Y\to X$为$T_0^{-1}$在$Y_0$上的延拓，且在$C$上为零，则$BT=I-\Pi_K$、$TB=I-\Pi_C$。若$J:X\to Y$紧，$S=T+J$满足$BS=I+K_1$、$SB=I+K_2$，两误差紧。由$u=BSu-K_1u$得
\[
 \|u\|_X\leq\|B\|\|Su\|_Y+\|K_1u\|_X.
\]
核的正交序列会与$K_1$的紧性矛盾。同样，若核正交补上没有$\|u\|_X\leq C\|Su\|_Y$，取范数1且$Su_j\to0$的序列；$K_1u_j$的收敛子列使$u_j$收敛到核中的单位向量，矛盾。因此像闭。对有界Hilbert伴随$S^*:Y\to X$用$B^*S^*=I+K_2^*$得其核有限维，所以$S$的余核有限维。这里紧算子的伴随仍紧：用有限$\epsilon$网逼近单位球像，并投影到网点张成的有限维空间，即以有限秩算子作范数逼近；取伴随后保持误差范数。

故$T+tJ$对全部$t\in[0,1]$均Fredholm，局部稳定性使其整数值指标局部常数，连通性使其常数。光滑零阶算子$H^1\to L^2$经紧嵌入分解，因而紧；联络改变只改变Dirac的零阶项。
\end{proof}

同一指标不表示核维数不变：题20.1给出核与余核同时跳跃的族。带边界时还须选择椭圆边界条件，非紧时须控制无穷远；本章不把它们包含在此定理中。


\Needspace{16\baselineskip}
\section{Atiyah--Singer定理及证明所需分析工具}
\subsection{特征形式提供的候选数}
第18章已定义$p_j$和$\operatorname{ch}$。以形式变量$x_j$表示Pontryagin根，规定$p=\prod_j(1+x_j^2)$，定义乘法序列
\begin{equation}\label{eq:ahat-expansion}
 \widehat A=\prod_j\frac{x_j/2}{\sinh(x_j/2)}
 =1-\frac{p_1}{24}+\frac{7p_1^2-4p_2}{5760}+\cdots.
\end{equation}
根仅为表达不变多项式的代数记号，不假定切丛全局分裂。函数为偶幂级数，关于$x_j^2$对称，故每个有限次数系数是$p_j$的多项式；第18章保证相应形式闭、类与联络无关。展开单因子为$1-x^2/24+7x^4/5760+\cdots$并相乘，直接给上述两项。本章不据此先验断言其积分为整数。

\Needspace{16\baselineskip}
\begin{theorem}[扭Spin指标公式；外引]\label{thm:as-external}
设$M^{2m}$紧致、定向、无边界且给定Spin结构，$E$为Hermite复丛并取酉联络。第19章手征约定下，
\begin{equation}\label{eq:as-spin}
 \operatorname{ind}D_E^+
 =\int_M[\widehat A(TM)\wedge\operatorname{ch}(E)]_{2m}.
\end{equation}
下标取顶次数。这里外引Atiyah--Singer定理的Spin特例，参见Dai讲义第3.3节定理3.3.2、3.3.4；\textbf{本书尚未证明此公式}。一般椭圆符号的拓扑指标、非紧Callias及带边APS公式不在本陈述内。
\end{theorem}

左侧已在定理\ref{thm:index-fredholm}定义并证明有限。右侧的闭性、联络无关性来自第18章；这些事实加上指标稳定性，并不能证明两侧相等。热核路线还需要热半群及光滑核、迹类与超迹配对、短时局部渐近和Getzler缩放。拓扑路线需要紧支$K$理论、Thom同构、Bott周期及解析指标与推前的比较。两条路线均有本书未建立的实质步骤，不能以“局部计算相同”替代。完整证明及一般整类比较列入补证任务。

\subsection{条件性推论与独立可算的例子}
二维$\widehat A$无正次数项，故承认\eqref{eq:as-spin}后
$\operatorname{ind}D_L^+=\int_Mc_1(L)$。
四维直接乘特征形式得
\begin{equation}\label{eq:as-four}
 \operatorname{ind}D_E^+=\int_M
 \left(\frac{c_1(E)^2}{2}-c_2(E)-\frac{\operatorname{rank}E}{24}p_1(TM)\right).
\end{equation}
这是外引定理的代数推论，不是另一份指标定理证明。若$E$平凡且联络平坦，右侧只剩$-r\int p_1/24$。若$M$没有给定Spin结构，不可直接套用未扭Spin公式；需先另行建立Spin$^c$或Dolbeault版本。

在单位圆球$S^2$上，第19章正标量曲率消失结论独立给两手征核零，所以指标零；$\widehat A$的二维部分也零，是相容检查。在平坦$T^2$的周期Spin结构中，两手征各有一个常旋量，指标仍零；其他符号结构两核均零，指标也零。核不同而指标相同并不矛盾。

若连通闭Spin流形具有$\mathrm{Scal}\geq0$且某点严格正，第19章给未扭$D$核零。再\emph{使用外引公式}，才能推出$\int_M\widehat A=0$。因此非零$\widehat A$数阻止这样的度量；扭丛曲率项有不定号，不能对任意$E$复用同一消失论证。题20.4核查这一边界。

\subsection{从指标到模空间尚差什么}
指标是线性化的虚维数来源。即使一个规范固定的线性算子Fredholm，非线性零点集也未必是同维数的光滑流形；仍需控制余核、稳定子、非线性映射的Banach光滑性及隐函数定理条件。若还要定义不变量，必须另外处理紧性、定向、扰动和度量改变。后文只证明明确写出的线性符号和变分恒等式。


\Needspace{16\baselineskip}
\section{Yang--Mills泛函与规范理论}
\subsection{联络仿射空间中的能量}
设$P\to M$为紧矩阵Lie群$G$的主丛，$\mathfrak g$上选正定Ad不变内积。对$G\subset U(r)$可取$\langle X,Y\rangle=-\operatorname{tr}(XY)$（视为实内积）。相伴丛$\operatorname{ad}P$承接它。第16章说明两个联络之差为$a\in\Omega^1(\operatorname{ad}P)$；局部
\[
 F_A=\dd A+A\wedge A,\qquad
 \dd_A\alpha=\dd\alpha+A\wedge\alpha-(-1)^k\alpha\wedge A,
 \quad\alpha\in\Omega^k(\operatorname{ad}P).
\]
用内积缩并Lie代数指标定义$L^2$配对与形式伴随$\dd_A^*$。在定向$n$维上它作用于$q$形式为$(-1)^{n(q+1)+1}*\dd_A*$；四维作用于二形式时为$-*\dd_A*$。定义无$1/2$因子的能量
\begin{equation}\label{eq:ym-energy}
 \mathcal Y(A)=\|F_A\|_{L^2}^2.
\end{equation}
该归一化使后面的SU2下界为$8\pi^2|k|$。换用半能量时所有能量常数也须减半。

\Needspace{17\baselineskip}
\begin{proposition}[曲率变分与Yang--Mills方程]\label{prop:ym-variation}
令$A_t=A+ta$，则
\begin{gather*}
 F_{A_t}=F_A+t\dd_Aa+t^2a\wedge a,\\
 \left.\frac{\dd}{\dd t}\right|_0\mathcal Y(A_t)
 =2(\dd_Aa,F_A)=2(a,\dd_A^*F_A).
\end{gather*}
故$A$对所有光滑联络变分临界当且仅当$\dd_A^*F_A=0$。此为Yang--Mills方程。
\end{proposition}
\begin{proof}
展开$\dd(A+ta)+(A+ta)\wedge(A+ta)$，一次项为$\dd a+A\wedge a+a\wedge A=\dd_Aa$，二次项为$a\wedge a$。对能量范数平方求导得第一配对。Ad不变性使局部协变分部积分相容，Stokes在无边界紧流形上给第二配对。若所有方向导数零，取$a=\dd_A^*F_A$即可使该截面范数平方为零；反向直接代入。
\end{proof}

第16章被动标架公式为$A^h=h^{-1}Ah+h^{-1}\dd h$、$F^h=h^{-1}Fh$，Ad不变性使能量不变。若用主动规范作用$g\cdot A=gAg^{-1}-(\dd g)g^{-1}$，它等于前式取$h=g^{-1}$；沿$g=e^{t\phi}$的切向量为$-\dd_A\phi$。全文在写主动作用时使用后一约定。

对$U(1)$写$A=i\alpha$，则$F=i\dd\alpha$，方程化为$\delta\dd\alpha=0$。在\emph{平凡线丛}的闭流形上，分部积分得$\|\dd\alpha\|^2=(\alpha,\delta\dd\alpha)=0$，故临界联络平坦；非平凡线丛没有全局$\alpha$，不能作同样推断。一般$U(1)$的曲率闭且余闭，因而为其上同调类的调和代表。题20.5给非零曲率实例。

\subsection{四维星分解及拓扑能量界}
现在$M$为定向四维。$*^2=1$作用于二形式，且星自伴，于是
\[
 \Lambda^2=\Lambda^+\oplus\Lambda^-,\qquad
 F^\pm=\tfrac12(F\pm *F),\qquad
 \|F\|^2=\|F^+\|^2+\|F^-\|^2.
\]
方程$F^-=0$称自对偶（SD），$F^+=0$称反自对偶（ASD）。改变底流形定向使星在二形式上变号，交换二者。

\Needspace{17\baselineskip}
\begin{theorem}\label{thm:ym-bound}
取$G=\mathrm{SU}(2)$及内积$-\operatorname{tr}(XY)$，记
\[
 k=\int_M c_2(E)=\frac1{8\pi^2}\int_M\operatorname{tr}(F_A\wedge F_A),
 \qquad E=P\times_{\mathrm{SU}(2)}\mathbb C^2.
\]
此数与联络无关，并有
\begin{equation}\label{eq:ym-bound}
 8\pi^2k=\|F_A^-\|^2-\|F_A^+\|^2,\qquad
 \mathcal Y(A)\geq8\pi^2|k|.
\end{equation}
当$k>0$时等号恰为ASD，$k<0$时恰为SD；$k=0$时等号恰为平坦。SD和ASD联络都满足Yang--Mills方程，并在固定丛的全部联络中达到绝对最小能量。
\end{theorem}
\begin{proof}
由于$\operatorname{tr}F=0$，第18章行列式展开给$c_2=\operatorname{tr}(F\wedge F)/(8\pi^2)$。联络无关性由传递公式和Stokes给出。此处先使用已证的实Chern--Weil数；其对一般丛的整数性仍依赖外部整类比较。

按$-\operatorname{tr}$的内积定义，
\[
 -\operatorname{tr}(F\wedge F)=\langle F,*F\rangle\mathrm{vol}
 =(|F^+|^2-|F^-|^2)\mathrm{vol}.
\]
积分得第一式，非负数之和不小于其差的绝对值得能量界，等号条件由较小一项为零得到。若$*F=\pm F$，Bianchi给$\dd_AF=0$，所以$\dd_A^*F=-*\dd_A*F=0$。固定丛上$k$不变，故达到该界就是全体联络中的最小值；不表示每个丛都存在达到下界的联络。
\end{proof}

\begin{figure}[htbp]
 \centering\DGDiagram{yang-mills-bound}

 \caption{固定$k>0$时的四维能量恒等式（原创示意）。横坐标为$\|F^+\|^2$，纵坐标为能量；直线来自$\mathcal Y=8\pi^2k+2\|F^+\|^2$。图不主张每个横坐标都能由给定丛的联络实现。}
\end{figure}

\subsection{共形不变与一个可算模型}
若$\widetilde g=e^{2f}g$，$q$形式的点积乘$e^{-2qf}$，体积乘$e^{nf}$。由星的定义$\alpha\wedge\widetilde *\beta=\langle\alpha,\beta\rangle_{\widetilde g}\mathrm{vol}_{\widetilde g}$得
\[
 \widetilde *|_{\Lambda^q}=e^{(n-2q)f}*.
\]
故四维二形式的星、SD/ASD方程和Yang--Mills能量均共形不变。它解释紧性分析为何会出现尺度问题，却不是起泡或可去奇点定理的证明。

在平坦单位体积$T^4$上，取线丛$L$的酉联络满足
$F_L=-2\pi i(a\,\dd x^{12}+b\,\dd x^{34})$，$a,b\in\mathbb Z$。
可在$\mathbb R^4$用$A_L=-2\pi i(ax^1\dd x^2+bx^3\dd x^4)$及整数平移的酉过渡函数下降，整数保证双重平移相容。记$\dd x^{ij}=\dd x^i\wedge\dd x^j$。在$E=L\oplus L^{-1}$取$F=\operatorname{diag}(F_L,-F_L)$，行列式联络平凡，故是SU2模型。直接计算
\[
 k=-2ab,\qquad \mathcal Y=8\pi^2(a^2+b^2).
\]
$b=-a$时ASD且$k=2a^2$；$b=a$时SD且$k=-2a^2$。曲率常系数且落在交换子代数中，故均为Yang--Mills；$a=1,b=0$则$k=0$而能量正，是临界却未达到下界的反例。这个可约模型不替代不可约BPST瞬子的全局构造。

\subsection{规范退化与椭圆复形的符号}
若$A$为ASD，$F_{A+ta}^+$的一次项是$\dd_A^+a=\pi_+\dd_Aa$。主动规范轨道的切向量为$-\dd_A\phi$，而
\[
 \dd_A^+\dd_A\phi=\pi_+[F_A,\phi]=[F_A^+,\phi]=0.
\]
所以出现复形
\begin{equation}\label{eq:asd-complex}
 \Omega^0(\operatorname{ad}P)\xrightarrow{\dd_A}
 \Omega^1(\operatorname{ad}P)\xrightarrow{\dd_A^+}
 \Omega^+(\operatorname{ad}P).
\end{equation}
在非ASD联络处，复合一般不为零，不能不加条件称它为复形。

\Needspace{17\baselineskip}
\begin{proposition}\label{prop:asd-symbol}
对每个实$\xi\ne0$，\eqref{eq:asd-complex}的符号列正合。规范固定算子
\[
 L_A=\dd_A^*\oplus\sqrt2\dd_A^+:
 \Omega^1(\operatorname{ad}P)\longrightarrow
 \Omega^0(\operatorname{ad}P)\oplus\Omega^+(\operatorname{ad}P)
\]
满足\eqref{eq:index-dirac-type}，故$H^1\to L^2$的延拓Fredholm。$\sqrt2$使符号恰为范数等距乘$|\xi|$；去掉它仍椭圆。
\end{proposition}
\begin{proof}
去掉不影响正合性的非零$i$因子，并把$\xi$归一化为正交余标架$e^1$。第一个映射$\phi\mapsto e^1\phi$单射。把$a=a_1e^1+\sum_{j=2}^4a_je^j$，则
\[
 \pi_+(e^1\wedge a)=\frac12\sum_{j=2}^4a_j(e^{1j}+*e^{1j}).
\]
括号中的三个自对偶形式正交且范数平方2，张成$\Lambda^+$。因此第二映射满射，其核恰为$e^1$方向，证明正合。$\dd_A^*$符号是$-i\iota_{\xi^\sharp}$，于是
\[
 |\iota_{\xi^\sharp}a|^2+2|\pi_+(\xi\wedge a)|^2
 =|\xi|^2|a|^2.
\]
两端实秩同为4乘$\dim G$，故单射符号也是满射，满足两个符号恒等式；复化后应用定理\ref{thm:index-fredholm}，或逐项使用其完全相同的实Hilbert证明。
\end{proof}

写复形的光滑上同调为$H_A^0,H_A^1,H_A^2$，分别为$\ker\dd_A$、$\ker\dd_A^+/\operatorname{im}\dd_A$和$\Omega^+/\operatorname{im}\dd_A^+$。此时
\begin{equation}\label{eq:asd-index-defect}
 \operatorname{ind}_{\mathbb R}L_A
 =\dim H_A^1-\dim H_A^0-\dim H_A^2.
\end{equation}
解析依据如下：$\dd_A^*\dd_A$的标量主符号为$|\xi|^2I$，其能量为$\|\dd_A\phi\|^2$；第14章的紧性、正交补估计与Hilbert变分证明给出其在核正交补上的解，光滑右端由同一差商证明光滑。因此每个$\ker\dd_A^+$的类可减去$\dd_A\phi$而满足$\dd_A^*a=0$，且代表唯一；这把$H_A^1$识别为$\ker L_A$。若$L_A^*(\phi,\eta)=\dd_A\phi+\sqrt2(\dd_A^+)^*\eta=0$，两项正交，因为$\dd_A^+\dd_A=0$，所以各为零。$L_A$的闭像还给$\dd_A^+$闭像：若$\dd_A^+a_j\to\eta$，先减去规范方向，则$L_Aa_j\to(0,\sqrt2\eta)$，再用闭像。由此余核正是$H_A^0\oplus H_A^2$。为使光滑商与Sobolev余核一致，对光滑$\eta$取其$L^2$到$\ker L_A^*$的投影；该有限维核光滑，余下右端在像内。求得$L_Aa$为此光滑右端后，$L_A^*L_A$的标量主部正则性使$a$光滑。

$H_A^0$描述无穷小稳定子，$H_A^2$是线性障碍空间；$H_A^2\ne0$不自动意味着每个形变受阻。即使两者零，仍需Sobolev乘积、局部规范切片和Banach隐函数定理才能从线性虚维数得到实际局部维数。本章尚未建立这些非线性工具，不由\eqref{eq:asd-index-defect}宣布模空间光滑、紧致或可定义Donaldson不变量。


\Needspace{16\baselineskip}
\section{Seiberg--Witten方程及适用条件}
\subsection{先给数据，再写方程}
四维Spin$^c$群定义为
\[
 \operatorname{Spin}^c(4)=(\operatorname{Spin}(4)\times U(1))/\{(1,1),(-1,-1)\}.
\]
在Spin表示上$[a,z]$作用为$z\rho(a)$，行列式角色为$z^2$。本节\emph{假定给定}到定向正交标架丛的Spin$^c$提升、秩二Hermite丛$W^\pm$、行列式线丛$L$以及相容Clifford作用。提升存在性、$c_1(L)\bmod2=w_2(TM)$与所有定向四维流形存在Spin$^c$结构的拓扑证明不在本章完成。

取$L$的酉联络$B$，与Levi--Civita联络合成旋量联络及$D_B^+:\Gamma(W^+)\to\Gamma(W^-)$。局部中心项为$B/2$，因为行列式角色是平方。对实二形式$\omega$约定$c(\omega)=\sum_{i<j}\omega_{ij}c_ic_j$。在第19章手征号下，$c$把$i\Lambda^+$实线性同构到$\operatorname{Herm}_0(W^+)$，把$i\Lambda^-$作用为零。具体地，以$e^{ij}=e^i\wedge e^j$记正向正交余标架，在Fock偶基$(1,u_1\wedge u_2)$上，三个自对偶基的像为
\[
 \begin{array}{c|c}
 i(e^{12}+e^{34})&\left(\begin{smallmatrix}2&0\\0&-2\end{smallmatrix}\right)\\
 i(e^{13}-e^{24})&\left(\begin{smallmatrix}0&-2i\\2i&0\end{smallmatrix}\right)\\
 i(e^{14}+e^{23})&\left(\begin{smallmatrix}0&-2\\-2&0\end{smallmatrix}\right)
 \end{array}
\]
它们构成无迹Hermite矩阵的实基，反自对偶三基作用为零；换标架由Spin等变性相容。这给出了同构的证明和尺度，而不是范数未定的“自然相等”。

定义无迹Hermite端同态
\[
 q(\psi)=\psi\psi^*-\tfrac12|\psi|^2I_{W^+}.
\]
本节固定以下未扰动Seiberg--Witten方程的归一化：
\begin{equation}\label{eq:sw-equations}
 D_B^+\psi=0,\qquad c(F_B^+)=q(\psi).
\end{equation}
有些资料把$1/2$或Clifford映射的尺度放到$q$中，比较时须同时换算平方公式；本书只用此处显式约定。

主动规范群$\mathcal G=C^\infty(M,U(1))$作用为
\begin{equation}\label{eq:sw-gauge}
 (B,\psi)\longmapsto(B-2u^{-1}\dd u,u\psi).
\end{equation}
因为旋量联络的中心项是$B/2$，直接对$u\psi$求协变导数可见$u^{-1}\dd u$两项抵消，故
\[
 D_{B-2u^{-1}\dd u}(u\psi)=uD_B\psi.
\]
又$\dd(u^{-1}\dd u)=0$，曲率不变，$q(u\psi)=q(\psi)$，所以方程规范不变。漏掉因子2便失去该协变性。

\Needspace{17\baselineskip}
\begin{proposition}[给定Spin$^c$数据的消失结论]\label{prop:sw-vanishing}
在本章紧致无边界假设下，若$\mathrm{Scal}\geq0$，\eqref{eq:sw-equations}的每个光滑解都有$\psi=0$且$F_B^+=0$。因此此时不存在$\psi\ne0$的未扰动解；并未断言全部解为空。
\end{proposition}
\begin{proof}
第19章Dirac平方的局部曲率缩并适用于Spin$^c$联络；中心曲率为$F_B/2$，故在$W^+$上
\[
 (D_B^+)^*D_B^+=\nabla_B^*\nabla_B+\tfrac14\mathrm{Scal}
 +\tfrac12 c(F_B^+).
\]
这是局部恒等式，推导只增加中心曲率，无需全局Spin分解。代入方程，与$\psi$配对并积分。由于$q(\psi)\psi=\tfrac12|\psi|^2\psi$，得到
\begin{equation}\label{eq:sw-integrated}
 0=\|\nabla_B\psi\|_2^2+\tfrac14\int_M\mathrm{Scal}|\psi|^2
 +\tfrac14\int_M|\psi|^4.
\end{equation}
每项非负，最后一项迫使$\psi=0$；第二条方程及$c:i\Lambda^+\to\operatorname{Herm}_0(W^+)$单射再给$F_B^+=0$。
\end{proof}

零旋量解称可约解，规范群中常数$U(1)$仍可固定它；若$M$连通且$\psi$某处非零，固定该对的规范变换由$\dd u=0$及$u\psi=\psi$迫使$u=1$。平坦$T^4$取平凡Spin$^c$数据及平坦$B$，$(B,0)$就是可约解，所以消失命题不能写成“解集为空”。

\subsection{研究对象及尚未建立的步骤}
模空间是\eqref{eq:sw-equations}的解集除以\eqref{eq:sw-gauge}。本章只定义这个集合，没有给它构造光滑流形结构。在完整外部理论中，规范固定形变算子的实指标为
\[
 d=\frac{\langle c_1(L)^2,[M]\rangle-(2\chi(M)+3\tau(M))}{4},
\]
其中$\chi$为Euler特征数、$\tau$为四维交叉形式的签名。\textbf{此虚维公式在本书仅作外引}；它依赖Spin$^c$指标、Hodge维数与签名公式的比较，不能由未知整类或尚未证明的签名定理跳步推出。Perutz第19.5节给出相应理论背景，本书未证明指标比较。

从虚维数到Seiberg--Witten不变量还需要Sobolev配置空间、椭圆规范切片、先验界、模规范紧性、横截扰动、同调定向和度量路径的配边。$b_2^+=1$时有腔室问题，不能无条件声称与度量无关；这里$b_2^+$为交叉形式正指数。外部理论对光滑结构的检测，不等于任何两个不同光滑结构都能被这些不变量区分。完整构造、壁穿越、紧性和与辛结构的关系留为后续专题。

\begin{figure}[htbp]
 \centering\DGDiagram{gauge-moduli-boundary}

 \caption{从方程到不变量的依赖（原创示意）。前两块中只有正文明确给出的变分、符号和消失命题已证明；箭头表示还需完成的分析与拓扑工作，不表示本书已有全部结果。}
\end{figure}


\Needspace{16\baselineskip}
\section{指标与规范场习题及详细解答}
以下8题均为依据本章定义设计的原创补充，不冒充周书或梁书原题。梁书相应主题的册次、题号及双页码尚未核实；外部讲义用于核对理论背景，未把讲义习题改号归为原创。20.3--20.4中使用指标公式的结论明确以外引定理为条件。

\Needspace{12\baselineskip}
\subsection*{习题20.1：核跳跃而指标不跳跃}
在单位长度圆周上取周期复函数，$P_a=\partial_x+ia:H^1\to L^2$，$a\in\mathbb R$。求核、余核、指标，并说明哪个结论需要闭像。
\begin{proof}[解答]
解$P_au=0$得$u(x)=Ce^{-iax}$，周期性要求$e^{-ia}=1$。故$a\in2\pi\mathbb Z$时核复维1，其他时候为零。形式伴随为$P_a^*=-\partial_x-ia=-P_a$，其核相同。主符号$i\xi$满足本章条件，定理\ref{thm:index-fredholm}给闭像并识别余核为伴随核，所以指标恒为零。若只解两个ODE而不证明像闭，就没有完成Fredholm判定。跨过$2\pi\mathbb Z$时两种缺陷同时跳跃，指标不变。
\end{proof}

\Needspace{12\baselineskip}
\subsection*{习题20.2：为什么非紧假设不能删除}
证明$\partial_x:H^1(\mathbb R)\to L^2(\mathbb R)$的核为零但像不闭。提示：缩放一个非零紧支光滑函数。
\begin{proof}[解答]
分布导数为零的$H^1$函数为几乎处处常数，$L^2(\mathbb R)$中的常数只能为零。取$\chi\in C_c^\infty(\mathbb R)$、$\|\chi\|_2=1$，令$u_R(x)=R^{-1/2}\chi(x/R)$，则
\[
 \|u_R\|_2=1,\qquad \|u_R'\|_2=R^{-1}\|\chi'\|_2.
\]
若像闭，单射的有界算子$\partial_x:H^1\to\operatorname{im}\partial_x$由有界逆定理给$\|u\|_{H^1}\leq C\|u'\|_2$，上述序列矛盾。因此像不闭，算子不Fredholm。其主符号处处可逆；失效的是无穷远的全局控制。这也展示了为何闭像不能从局部椭圆性直接推出。
\end{proof}

\Needspace{12\baselineskip}
\subsection*{习题20.3：四维扭指标的次数核对}
设$M^4$闭且有给定Spin结构。对秩$r$复丛$E$，从$\widehat A$和$\operatorname{ch}$展开推出\eqref{eq:as-four}。当$E=L\oplus L^{-1}$时把结果写成$c_1(L)$。哪些步骤依赖外引？
\begin{proof}[解答]
四维只保留微分形式次数0、2、4：
\[
 \widehat A=1-\frac{p_1}{24},\qquad
 \operatorname{ch}(E)=r+c_1+\frac12(c_1^2-2c_2).
\]
因$\widehat A$没有二次形式项，乘积的四次部分恰为$\frac12c_1^2-c_2-rp_1/24$。对$E=L\oplus L^{-1}$，Whitney给$c_1(E)=0$、$c_2(E)=-c_1(L)^2$、$r=2$，所以
\[
 \operatorname{ind}D_E^+=\int_M\left(c_1(L)^2-\frac{p_1(TM)}{12}\right).
\]
特征形式代数由第18章和幂级数展开证明；把其积分认作解析指标这一步依赖定理\ref{thm:as-external}，本题没有补齐该定理。
\end{proof}

\Needspace{12\baselineskip}
\subsection*{习题20.4：未扭消失不能推广到任意扭丛}
设闭Spin曲面取线丛$L$，$\int c_1(L)=d\ne0$。承认外引指标公式，证明$\dim\ker D_L^++\dim\ker D_L^-\geq|d|$。这为何与正标量曲率下的未扭消失不冲突？
\begin{proof}[解答]
二维公式给$\dim\ker D_L^+-\dim\ker D_L^-=d$。两个非负整数之和不小于其差的绝对值，得到所需下界。未扭平方只有粗Laplace与$\mathrm{Scal}/4$；扭平方还含$\sum_{i<j}c_ic_jF^L_{ij}$，它没有统一非负号。故即使底流形具有正标量曲率，未扭的积分非负性也不能迫使扭核为零。例子可取具有第19章Spin结构的$S^2$及第16章Hopf线丛的非零次幂；此处核的具体维数未由不等式确定。
\end{proof}

\Needspace{12\baselineskip}
\subsection*{习题20.5：环面上的非零Yang--Mills曲率}
补全正文$T^4$上线丛的下降构造，计算$L\oplus L^{-1}$的能量与$k$。分别判断参数$(a,b)$取$(1,-1)$、$(1,1)$和$(1,0)$时的三例。
\begin{proof}[解答]
在$\mathbb R^4$取$A=-2\pi i(ax^1\dd x^2+bx^3\dd x^4)$。平移$x^1\mapsto x^1+1$时$A$减去$2\pi ia\dd x^2$，可用被动酉换标架$h_1=e^{-2\pi iax^2}$实现；$x^3$方向同理用$h_3=e^{-2\pi ibx^4}$，其他方向用1。由于$a,b$为整数，$h_1$在$x^2$的整数平移下不变，$h_3$在$x^4$的整数平移下不变；两个非平凡过渡又交换，所以所有格平移的余循环相容。得到线丛与联络，而非平凡线丛上的全局一形式。

令$\alpha=a\dd x^{12}+b\dd x^{34}$，则$F_L=-2\pi i\alpha$、$c_1(L)=\alpha$。Whitney给$c_2(L\oplus L^{-1})=-\alpha^2$，$\int\alpha^2=2ab$，故$k=-2ab$。矩阵迹内积给$|F|^2=8\pi^2(a^2+b^2)$，体积1得到能量。$*(\dd x^{12})=\dd x^{34}$：第一例ASD，$k=2$、能量$16\pi^2$；第二例SD，$k=-2$、能量相同；第三例既非SD也非ASD，$k=0$、能量$8\pi^2$。三者曲率均平行且与联络交换，故都是Yang--Mills。第三例显示临界点不必达到拓扑下界。
\end{proof}

\Needspace{12\baselineskip}
\subsection*{习题20.6：四维能量和定向}
证明四维二形式能量对共形变换不变。反转定向后$k$、$\mathcal Y$与ASD标签怎样改变？为什么不能只改$k$而保留原来的正负分解？
\begin{proof}[解答]
对$\widetilde g=e^{2f}g$有$|F|_{\widetilde g}^2=e^{-4f}|F|_g^2$、$\mathrm{vol}_{\widetilde g}=e^{4f}\mathrm{vol}_g$，乘积不变。反转定向使顶次形式积分换号，故$k\mapsto-k$；能量按正的体积密度积分，不变。星由定向定义，在二形式上变成$-* $，所以$F^+$与$F^-$交换，原ASD成为SD。此时
$\|F^-\|^2-\|F^+\|^2$也换号，与$8\pi^2k$同步。若只改$k$，便破坏同一恒等式的定向约定。
\end{proof}

\Needspace{12\baselineskip}
\subsection*{习题20.7：规范方向与缺少的符号秩}
令$\xi=e^1$。计算$\pi_+(\xi\wedge a)$的核，说明$\dd_A^+$单独不能成为同秩椭圆算子，再验证加上$\dd_A^*$如何恢复可逆符号。此论证是否证明非线性模空间光滑？
\begin{proof}[解答]
写$a=\sum a_je^j$，得到$\pi_+(e^1\wedge a)=\sum_{j=2}^4a_j(e^{1j}+*e^{1j})/2$。三个括号线性独立，故核恰为$a=a_1e^1$，是第一个符号$\xi\phi$的像。原映射秩从4到3且有一维核，不能满足本章可逆符号的定义。加上$-\iota_{\xi^\sharp}$提取$-a_1$，而$\sqrt2\pi_+(\xi\wedge a)$保留其余三分量的范数，得到$|\sigma_{L_A}(\xi)a|^2=|a|^2$。恢复的是规范固定的线性椭圆性；非线性切片、余核和稳定子问题仍在，不能据此认定模空间光滑。
\end{proof}

\Needspace{12\baselineskip}
\subsection*{习题20.8：Spin$^c$因子与四次正项}
用局部中心项$B/2$验证\eqref{eq:sw-gauge}；对$\psi\in\mathbb C^2$求$q(\psi)$的特征值、Hilbert--Schmidt范数和$q(\psi)\psi$。解释$\mathrm{Scal}\geq0$时为何只能推出零旋量而不能一概说无解。
\begin{proof}[解答]
设$u=e^{i\theta}$局部成立，则$B'=B-2i\dd\theta$。展开
\[
 (\dd+B'/2)(u\psi)
 =u(\dd+B/2)\psi+(\dd u-ui\dd\theta)\psi
 =u(\dd+B/2)\psi.
\]
Levi--Civita部分与标量$u$交换，因此Dirac也协变。若错误地只减$u^{-1}\dd u$，会余下$(\dd u)\psi/2$。

令$s=|\psi|^2$。非零$\psi$张成方向上$\psi\psi^*$特征值为$s$，正交方向为0，所以$q$特征值为$s/2,-s/2$，
\[
 |q(\psi)|_{\mathrm{HS}}^2=s^2/2,\qquad q(\psi)\psi=(s/2)\psi.
\]
当$\psi=0$这些式子仍成立。平方公式中的$1/2$与$q$作用的$1/2$相乘，正是积分恒等式中的$|\psi|^4/4$。非负标量曲率迫使$\psi=0$，随后要求$F_B^+=0$；若存在这样的行列式联络，它仍给可约解。平坦$T^4$的平凡$B$就是例子。
\end{proof}

